% % !TeX root = ../main.tex

% \chapter{补充材料}
% \label{chap:complementary}

% 本附录给出与正文对应的实现补充，重点说明两条主线中的关键数据结构与核心流程。为避免与工程代码逐字对应，以下均采用抽象命名，只保留有助于复现与理解的必要信息。

% \section{全局布线侧补充说明}
% \label{sec:appendix-routing-impl}

% \subsection{关键数据结构（抽象）}

% \begin{table}[!htbp]
% 	\centering
% 	\small
% 	\caption{全局布线侧关键数据结构（抽象表示）}
% 	\label{tab:appendix-routing-data-structures}
% 	\begin{tabular}{p{0.25\textwidth}p{0.30\textwidth}p{0.35\textwidth}}
% 		\hline
% 		抽象对象 & 关键字段 & 作用 \\
% 		\hline
% 		网络状态单元 & 网络编号、引脚集合、当前树结构 & 维护单网在不同阶段的布线状态。 \\
% 		栅格资源单元 & 容量、需求、边代价 & 反映通道资源使用情况并驱动拥塞判断。 \\
% 		候选树集合 & 候选拓扑、候选评分、选中索引 & 支撑动态更新中的树切换与重选。 \\
% 		路径决策图 & 节点层区间、候选边、回溯指针 & 承载层分配与路径回溯过程。 \\
% 		统计结果单元 & 线长、过孔数、运行时间 & 用于阶段评估与跨轮比较。 \\
% 		\hline
% 	\end{tabular}
% \end{table}

% \subsection{核心流程伪代码}

% \begin{algorithm}[!htbp]
% 	\KwIn{网络集合 $\mathcal{N}$，栅格资源图 $G$，动态轮数 $R$}
% 	\KwOut{最终布线树集合}
% 	\ForEach{$n \in \mathcal{N}$}{
% 		构建初始拓扑并完成一次模式布线与层分配\;
% 		提交到资源图 $G$\;
% 	}
% 	\For{$r\leftarrow1$ \KwTo $R$}{
% 		提取拥塞网络集合 $\mathcal{N}_{of}$；若为空则提前结束\;
% 		\ForEach{$n\in\mathcal{N}_{of}$}{
% 			撤销旧结果并更新候选拓扑\;
% 			重新执行模式布线与层分配\;
% 			提交新结果\;
% 		}
% 	}
% 	对仍拥塞网络执行细粒度搜索修复\;
% 	输出最终统计指标\;
% 	\caption{拥塞感知初始布线主流程（抽象）}
% 	\label{alg:appendix-routing-main}
% \end{algorithm}

% 上述流程与第三章一致：先保证高效初始解，再通过拥塞反馈进行增量修正，最后对剩余困难网络执行补充搜索。

% \section{布局侧补充说明}
% \label{sec:appendix-placement-impl}

% \subsection{关键数据结构（抽象）}

% \begin{table}[!htbp]
% 	\centering
% 	\small
% 	\caption{布局细粒度探索关键数据结构（抽象表示）}
% 	\label{tab:appendix-placement-data-structures}
% 	\begin{tabular}{p{0.25\textwidth}p{0.30\textwidth}p{0.35\textwidth}}
% 		\hline
% 		抽象对象 & 关键字段 & 作用 \\
% 		\hline
% 		约束描述单元 & 区域边界、约束强度、对象集合 & 表示可被优化的布局引导约束。 \\
% 		几何扰动单元 & 平移、伸缩、边界微调、类型切换 & 定义单步邻域操作。 \\
% 		评估结果单元 & 线长、过孔、运行时间等指标 & 作为外层搜索的反馈信号。 \\
% 		搜索状态单元 & 温度、当前最优值、停滞轮数 & 控制多段退火的探索与收敛。 \\
% 		\hline
% 	\end{tabular}
% \end{table}

% \subsection{多段退火优化伪代码}

% \begin{algorithm}[!htbp]
% 	\KwIn{初始约束 $C_0$，每轮迭代上限 $K$，最大停滞轮数 $S_{max}$}
% 	\KwOut{最优约束 $C_{best}$}
% 	$C_{best}\leftarrow C_0$，$s\leftarrow 0$\;
% 	\While{$s<S_{max}$}{
% 		以当前最优解作为新一轮初值\;
% 		初始化温度 $T$ 并执行 $K$ 次退火迭代\;
% 		每次迭代执行：生成邻域扰动 $\rightarrow$ 工具评估 $\rightarrow$ 按 Metropolis 准则接受/拒绝\;
% 		\If{本轮最优优于 $C_{best}$}{
% 			更新 $C_{best}$，$s\leftarrow0$\;
% 		}
% 		\Else{
% 			$s\leftarrow s+1$\;
% 		}
% 	}
% 	返回 $C_{best}$\;
% 	\caption{基于多段退火的外层细粒度优化（抽象）}
% 	\label{alg:appendix-placement-sa}
% \end{algorithm}

% \section{附录说明}

% 本附录仅保留复现所需的核心结构与流程。更细粒度的工程实现细节（例如具体脚本接口、日志字段和中间文件组织）可在开源实现中进一步查阅。

% \endinput



